\documentclass[dvipsnames,usenames]{beamer}
\usepackage{etex}
%\input{page}
%\usepackage{fullpage}	% Does about the same as the \input{page}...
% ^ part of the preprint package...

\input{symbolmac}
\usepackage{amsmath}
\usepackage{amssymb}  	% provides Rbb, etc.
\usepackage{latexsym}  	% provides \Box
%\usepackage{theorem} \input{proof}
%\pagestyle{empty}
\usepackage{graphicx} 	% for inserting graphics.
%\usepackage{booktabs}	% for fancy tabular environments and stuff
%\usepackage{subfig}	% for subfigures... 
%(This breaks Beamer if the patch isn't applied!!)
%usepackage{hyperref}

\definecolor{usuBlue}{RGB}{11,51,94}
\usetheme[not@ku={Utah State University},
          basecolor=usuBlue,
          topbarcolor=usuBlue,
          headstyle=institute]{Frederiksberg}
\institute[ECE Department]{\hfill Department of Electrical and
  Computer Engineering}
%\logo{\includegraphics[width=.1\textwidth]{horizontal-stacked-logo-blue}}
\logo{\includegraphics[width=.1\textwidth]{./ece_logo}}

%% ******** For code: ******************
%\usepackage{listings}
%\usepackage[usenames,dvipsnames]{color}
%% ******** Define the language: *******
%\lstdefinelanguage{matlabfloz}{%
%	alsoletter={...},%
%	morekeywords={% 						% keywords
%	break,case,catch,continue,elseif,else,end,for,function,global,%
%	if,otherwise,persistent,return,switch,try,while,...,
%        classdef,properties,methods},%
%	comment=[l]\%,% 						% comments
%	morecomment=[l]...,% 					% comments
%	morestring=[m]',%   						% strings 
%}[keywords,comments,strings]%
%\lstdefinestyle{matlab}{language=matlabfloz,							
%		keywordstyle=\color[rgb]{0,0,1},					
%		commentstyle=\color[rgb]{0.133,0.545,0.133},	% comments
%		stringstyle=\color[rgb]{0.627,0.126,0.941},	% strings
%		numbersep=3mm, numbers=left, numberstyle=\tiny% number style
%}
%% ******** Set the style: **************
%\lstset{style=matlab}

%\newcommand{\matlab}{{\sc Matlab }} % Or 
\renewcommand{\matlab}{{\sc Matlab }} 

% \newcommand{\homedir}{/home/aepound}
% \newcommand{\figdir}{\homedir/figs}  % picture stuff: ps,fig,tex
%\newtheorem{theorem}{Theorem}
%\newtheorem{corollary}{Corollary}
%\newtheorem{lemma}{Lemma}

\title{Back-mounted Jetpack \\
Thrust Vector Control\\}
\subtitle{
ECE 6320, Fall 2014
}
\author{Andrew Pound, Dean Lanier}
\date{\today}



\begin{document}


\KUtitlepageULimage[height=.157\textheight]%
        {./horizontal-stacked-logo-blue} 
% {horizontal_wordmark_blue}
\KUtitlepageLRimage{./eng-innov2}

%\maketitle
\frame[plain]{\titlepage}

\begin{frame}
\frametitle{Table of contents}
\tableofcontents
\end{frame}

\section{Background}

\begin{frame}
\frametitle{Introduction}  
Project is based on the 4MM Projec that has been developed 
at ASU.

\vspace{1em}
\href{https://www.youtube.com/watch?v=vJSbd8ttLoU}{YouTube Video}
\vspace{1em}
\begin{figure}
  \centering
  \includegraphics[width=.5\linewidth]{../figs/jetpack}
  \caption{Backpack Jetpack}
\end{figure}

\end{frame}

\begin{frame}
  \frametitle{Main Difficulties of Project}
  \begin{itemize}
  \item Equations of motion for human running are not
    available/tractable (at least not to us!)
  \vspace{3em}
  \item We don't have a rocket, and thus don't have  a model for how
    it reacts.
  \end{itemize}
\end{frame}


\begin{frame}
\frametitle{Motion Capture Data}  

We needed to be able to see what the back of a person that was running
was like.  

\vspace{1em}
We were able to find a few different research databases with Motion
Capture information of a person running.

\vspace{1em}
Motion Capture databases:
\begin{itemize}
\item ACCAD at OSU
\item MoCapClub
\item CMU MoCap database
\end{itemize}
\end{frame}

\section{Modelling and Controlling}
\begin{frame}
\frametitle{Assumptions and Simplifications}

\begin{itemize}
\item Flat ground
\vspace{1em}
\item Simplified model
\begin{itemize}
\item 2nd order
\item No coupling between angular coordinates
\end{itemize}
\vspace{1em}
\item Placement of sensors
\vspace{1em}
\item Simplification of body model
\begin{itemize}
\item Runner exhibits sinusoidal movements
\item Back plane based at hips
\end{itemize}
\end{itemize}

\end{frame}


\begin{frame}
\frametitle{Controller}
%Pole placement feedback control.
We used a pole placement state feedback controller to provide the
control to the rocket thrust.

\begin{itemize}
\item Used a 2nd order system for the rocket model
%\vspace{1em}
%\item 
\end{itemize}
\begin{figure}
  \centering
  \includegraphics[width=.7\linewidth]{../figs/sim}
  \caption{Simulink model}
\end{figure}

\end{frame}

\section{Results}
\begin{frame}
  \frametitle{Results}
  Simulation time!
\end{frame}

\section{Conclusions \& Future Work}

\begin{frame}
\frametitle{Future Work}
Conclusions
\begin{itemize}
\item Controlling the thrust vector works! 
%\vspace{1em}
%\item 
\end{itemize}

\vspace{1em}
Future Work
\begin{itemize}
\item Calculation of efficiency gain
\vspace{1em}
\item Better specification of rocket model
\vspace{1em}
\item More accuracy of back movements
\vspace{1em}
\item Explore different state-space representations
\end{itemize}
\end{frame}




%\begin{abstract}
%\end{abstract}




% \bibliographystyle{ieeetr}
% \newcommand{\bibdir}{~}
% \bibliography{\bibdir/~}
\end{document}
